% ===================== 第 3 章 =====================
\section{系统架构总览}

\subsection{模块化单体：一个进程内的清晰边界}

Poliscope 采用模块化单体加异步 Worker，首版不做微服务拆分。这是刻意的范围控制：MVP 只需要证明「这套组织结构比普通研究智能体更能发现高价值盲点」，拆分服务对此没有贡献，反而会引入分布式一致性问题。模块边界在代码层强制执行——跨模块只允许通过明确的 Pydantic 契约与应用服务通信，不能直接操作对方的数据表。

\begin{table}[ht]
\centering\small
\caption{14 个内部包的职责划分（\code{packages/}）}
\label{tab:modules}
\begin{tabularx}{\linewidth}{@{}L{3.1cm}X@{}}
\toprule
\rowcolor{accent}
\tblhead 模块 & \tblhead 职责 \\
\midrule
\code{kernel}      & 冻结契约、数据库会话、配置与权限常量；被所有模块依赖，自身不依赖业务模块 \\
\code{council}     & 七席位规格、结构化动作、八阶段注册表、席位与网关的桥接 \\
\code{epistemo}    & 编排器、显式状态机、预算追踪、停止条件、断点恢复 \\
\code{memory}      & MemoBrain 适配器、席位私有记忆、过程图与折叠 \\
\code{evidence}    & 事件账本、证据门、图投影器、谱系、独立性、辩证折叠 \\
\code{papers}      & 查询规划、多源候选发现、规范化、全文解析、发现抽取 \\
\code{models}      & 统一模型网关、成本审计、录制回放、免费体验额度 \\
\code{tools}       & 统一工具网关与四个学术数据源适配器 \\
\code{research}    & 研究契约、主张原子化、任务仓储与应用服务 \\
\code{reports}     & Research Brief 与最终论文组装、Markdown/JSON 双格式、导出脱敏 \\
\code{evaluation}  & ForesightBlindspot 语料、基线/消融接线、语义裁判、标注一致性 \\
\code{accounts}    & 注册登录、邮箱验证码、PBKDF2 口令、bearer 会话 \\
\code{knowledge}   & 个人知识库：多格式解析、全文检索、跨任务复用 \\
\code{skills}      & 从 GitHub 安装研究者技能（SKILL.md），按任务注入议会提示词 \\
\bottomrule
\end{tabularx}
\end{table}

应用层四个入口共享同一套后端契约：\code{apps/api}（FastAPI HTTP 层）、\code{apps/worker}（任务认领与议会执行）、\code{apps/web}（React 证据工作台）、\code{apps/cli}（纯 HTTP 客户端）。CLI 被硬性禁止 import 任何 \code{packages}——第二条直连研究服务的路径，就是一条绕过证据门的路径。

\subsection{一次研究的生命周期}

\begin{enumerate}[leftmargin=1.6em, itemsep=3pt]
\item \textbf{建任务。}研究者提交问题、范围、预算与可选的知识库、PDF、技能；系统原子化出建议主张，任务停在 \code{AWAITING\_CLAIM\_CONFIRMATION}，研究者不确认就不开跑。
\item \textbf{确认主张。}研究者勾选要调查的原子主张后任务入队；被放弃的主张标记为 \code{DISCARDED} 而非删除。
\item \textbf{Worker 认领。}用 \code{SELECT ... FOR UPDATE SKIP LOCKED} 从 PostgreSQL 队列认领，多个 Worker 并发安全；崩溃遗留的 \code{RUNNING} 由 \code{recover\_stale\_running} 回收。
\item \textbf{议会执行。}七席乘八阶段在一次数据库事务内跑完，每席经统一网关与模型交互，结构化动作追加进事件账本。
\item \textbf{图投影。}投影器以第二个数据库身份、第二个事务消费账本事件，过证据门后写入证据图。
\item \textbf{报告。}终态时从证据图、账本与任务行组装 Research Brief 和最终论文，报告阶段不调用模型产生新判断。
\end{enumerate}

\subsection{架构全景}

\begin{figure}[ht]
\centering
\resizebox{\linewidth}{!}{%
\begin{tikzpicture}[node distance=2.6mm and 6mm]
  \node[ndd, minimum width=30mm] (web) {Web 工作台 / CLI};
  \node[nd, below=of web, minimum width=30mm] (api) {FastAPI 应用层（apps/api）};
  \node[nda, below=of api, minimum width=30mm] (worker) {Worker：七席 × 八阶段议会};
  % left rail
  \node[nd, left=11mm of worker, minimum width=24mm] (mg) {模型网关\newline\tiny 分层模型 · 审计};
  \node[nd, below=of mg, minimum width=24mm] (tg) {工具网关\newline\tiny 四大学术源};
  % center chain
  \node[nd, below=of worker, minimum width=30mm] (ledger) {科学事件账本（追加写 · 幂等）};
  \node[nd, below=of ledger, minimum width=30mm] (gate) {六阶段证据门};
  \node[ndd, below=of gate, minimum width=30mm] (proj) {Graph Projector（唯一写入者）};
  \node[nda, below=of proj, minimum width=30mm] (graph) {证据图（10 类节点 / 12 类边）};
  % right rail
  \node[nd, right=11mm of ledger, minimum width=24mm] (proc) {过程图 + 三层记忆\newline\tiny Flush/Fold/Recall};
  \node[nd, below=of proc, minimum width=24mm] (bs) {盲点悬赏\newline\tiny 回流再调查};
  \draw[flow] (web)--(api); \draw[flow] (api)--(worker);
  \draw[flow] (worker)--(ledger); \draw[flow] (ledger)--(gate);
  \draw[flow] (gate)--(proj); \draw[flow] (proj)--(graph);
  \draw[flowg] (worker)--(mg); \draw[flowg] (worker)--(tg);
  \draw[flowg] (worker)--(proc);
  \draw[flow] (graph)--(bs); \draw[flowg] (bs.north) to[bend left=18] (worker.east);
\end{tikzpicture}}
\caption{系统架构：应用、议会、账本、证据门、投影器构成主干；模型/工具网关是唯一外部出口；过程图与证据图分居两侧，盲点回流驱动增量调查。}
\label{fig:arch}
\end{figure}

\subsection{两个事务、两个数据库身份}

「投影器是证据图唯一写入者」如果只写在文档里，就只是愿望。Poliscope 在 PostgreSQL 角色层把它变成硬事实：

\begin{itemize}[leftmargin=1.6em, itemsep=3pt]
\item 应用身份 \code{poliscope\_app} 可以追加账本、读取证据图，但对证据图表没有 \code{INSERT/UPDATE/TRUNCATE} 权限；
\item 投影器身份 \code{poliscope\_projector} 可以写图、推进事件状态，但不能改写既有事件 payload，且没有 \code{DELETE}（会话删除是唯一的迁移级例外，见第 9.7 节）；
\item 两个身份都不能 \code{TRUNCATE}。
\end{itemize}

Worker 因此把一次运行拆成两个事务：议会事务以应用身份提交动作，投影事务以投影器身份落图。集成测试 \code{test\_graph\_write\_isolation.py} 逐条断言应用身份写图会被数据库直接拒绝——这不是代码自律，是数据库事实。

\subsection{为什么队列建在 PostgreSQL 而不是 Redis}

任务认领与状态更新必须在同一个事务里完成，Worker 崩溃时行锁随事务自动释放，不会留下「谁在跑」的悬空状态。引入 Redis 等于多一个可能与数据库不一致的真相来源，而首版没有任何组件需要 Redis 做缓存或广播。少一个服务，就少一类故障与一份运维成本。
